\paragraph{The reduced soft endpoint.}
For an explicit closed endpoint expression, apply the finite
primitive operator $\mathfrak T$ just defined to the four reduced
kernels and the logarithm lists $\mathcal A_\beta,\mathcal A_1,\mathcal A_k$. The notation $\mathfrak T[R;L_aL_b]$ means the finite
sum over their weights after applying
$G(a)G(b)=G(a,b)+G(b,a)$.  Put
\begin{equation}
 \begin{aligned}
 \widetilde\Psi(z)={}&
 \mathfrak T[\widetilde R_0;L_\beta](z)
 +\mathfrak T[\widetilde R_1;L_\beta L_1](z)\\
 &+\mathfrak T[\widetilde R_k;L_\beta L_k](z)
 +\mathfrak T[\widetilde R_f;L_k](z).
 \end{aligned}
\end{equation}
Each primitive is normalized to vanish at $z=0$.  The simplified
answer is
\begin{equation}
 \boxed{\mathcal C_K(s)=
 \lim_{z\to Z^-}\widetilde\Psi(z)
 -\frac{\operatorname{artanh}\beta}{\beta}\mathcal C_0(s),
 \qquad Z=\sqrt s+\sqrt X-1.}
 \label{eq:comp-ck-simplified-endpoint}
\end{equation}
The combined endpoint is finite.  All terms are GPLs of weight at
most three with algebraic endpoint coefficients.  There are two further exact residue cancellations.  If
$P_1$ denotes the primitive of the higher-order pole part of
$\widetilde R_1$ at $w=0$, then
\begin{equation}
 \operatorname*{Res}_{w=0}\widetilde R_1=0,\qquad
 \operatorname*{Res}_{w=0}\left[
 \widetilde R_0-P_1\frac{\dd L_1}{\dd w}\right]=0.
\end{equation}
The first identity removes 32 weight-three words introduced by a
vanishing simple-pole coefficient.  The second removes four
weight-two words.  They follow by substituting the displayed rational
kernels, so no small numerical coefficient is discarded.
After these cancellations, the finite rule produces 190 terms instead
of 306, and removes both imaginary letters.  The same physical interval fixes all logarithms.
